\documentclass[letterpaper]{article}

\usepackage[margin=1.5cm]{geometry}
\usepackage{amsmath,amsfonts}
\usepackage{url}
\usepackage{bm}

\newcommand{\Real}{\mathbb{R}}
\newcommand{\tr}{^\top}
\newcommand{\opt}{^\star}

\setlength{\parskip}{3mm plus1mm minus1mm}
\setlength{\parindent}{0mm}

\title{Problem Set 5}
\author{Foundations of Machine Learning}
\date{}

\begin{document}
\maketitle

\begin{itemize}
\item \textbf{Submission}: Turn in as a \texttt{PDF} and a \textbf{separate source code file} on \url{http://mycourses.unh.edu}. Scanned hand-written notes are encouraged.
\item \textbf{Questions}: See the class website for a link to office hours
\end{itemize}

% Covers: Vector/quadratic functions using linear algebra (CO A.3--A.5) and
% Multivariate calculus basics: gradient and Hessian (CO A.4, MML 5.2).
% Boyd's appendix has no exercises, so the calculus problems are drawn from MML ch.~5.

\textbf{Relevant reading}: CO A.3--A.5; MML 5.1--5.2. Problems marked \textbf{[code]} should be submitted as a program in a language of your choice (Julia or Python recommended). The cited exercise numbers are from MML (Deisenroth et al., 2021).

\begin{enumerate}
\item MML, Chapter 5 exercises: any three on partial derivatives, gradients, and the Hessian (e.g.\ Exercises 5.1--5.4).
\item For the quadratic form $f(x) = x\tr A x$ with $A$ symmetric, derive the gradient $\nabla f$ and the Hessian $\nabla^2 f$ by hand, and explain how the definiteness of $A$ (via its eigenvalues) determines whether $f$ has a minimum, a maximum, or a saddle at the origin.
\item \textbf{[code]} For $f(x,y) = x^2 + 2xy + 3y^2$, compute the gradient and Hessian by hand, then verify both numerically by finite differences at the point $(1,1)$. Plot the contours of $f$.
\item \textbf{Time spent.} Report approximately how many hours you spent completing this problem set (including the reading).
\end{enumerate}

\end{document}
